% page 514 du fac-simile (image p-513 du scan)
% bloc 198.1 x 207.4 mm ; marge gauche 3.0 mm
\pgc{-17.780mm}{0.000mm}{
\l{}
\l{}
\l{}
\l{}
\l{}
\l{\cel{18}506.}
\l{}
\l{\cel{15}\sou{ruptar}. (\sou{trans.}) Pecigar, pecetigar, per shoko, frapo, preso}
\l{\cel{18}tre forta. - Anke metafore. - EFIS.}
\l{}
\l{\cel{15}\sou{ruro}. Extensajo de sulo, situita exter la urbi, quan okupas}
\l{\cel{18}la agri, la prati, e c. - EFIS.}
\l{}
\l{\cel{15}\sou{ruaDe}. Bendo de stofo, de tulo, e c., plisita fomme ii aDom}
\l{\cel{18}kupeti, di godroni, kelke simila a la alveoli di abeluyo,}
\l{\cel{18}ed uzata por ornar la vesti hominala. - DEFR.}
\l{}
\l{\cel{15}\sou{rusko.(bot.}) Genero de planti ek la familio \textquotedbl{}liliazei\textquotedbl{}, m mo}
\l{\cel{18}roti eempre verda, simila ad la asparagi, kun folii rmihmnna}
\l{\cel{18}e rameti aplastita forme di folii verda e harda ledratre}
\l{\cel{18}di qui la surfaco portas flori qui donas beri globatra.}
\l{\cel{18}- L.\cel{1}\sou{ruscus}.-ISL.}
\l{}
\l{\cel{15}\sou{rusto}. Fer-oxido, redatra. -\cel{1}\sou{Simb. kem}.: Fe2 O3 e Fe3 O4. DE.}
\l{}
\l{\cel{15}\sou{rustika}. Qua apartenas a la kozi di ruro.\cel{2}Dicesas anke pri}
\l{\cel{18}la manieri di la rurani - pri la planti qui rezistas la}
\l{\cel{18}domaji produktita da la mala vetero quale la planti sovaja}
\l{\cel{18}- pri benko, kozo ek petro kruda od ek ligno ne-laborita.}
\l{\cel{18}- DEFIS.}
\l{}
\l{\cel{15}\sou{ruto}. (\sou{bot.}) Planto medicinala, ek la familio \textquotedbl{}rutacei\textquotedbl{} di}
\l{\cel{18}qui lu esas la tipo, qua saporas akre. - L.\cel{1}\sou{ruta}. - DFIRS.}
\l{}
\l{\cel{15}\sou{ruteno}. (\sou{kemio)}. Metalo ek la grupo di platino, la maxim skar-}
\l{\cel{18}sa ek omni, deskovrita da Glaus (en la yaro 1843) en la}
\l{\cel{18}osmio di iridio. -\cel{1}\sou{Simbolo kemiala}\cel{1}: Ru.}
\l{}
\l{\cel{15}\sou{rutino,}\cel{2}Kustumo su-inkrustinta agar o facar ulo seg1s ta o}
\l{\cel{18}maniero. -\cel{1}\sou{Metaf.}:Kustumo a qua onu su konformas mekanis-}
\l{\cel{18}matre, sen-reflekte. - DEFIRS.}
\l{}
\l{\cel{15}\sou{ruzar.}\cel{1}(\sou{netrans., per)}\cel{2}Uzar, por trompar, moyeni qui igas}
\l{\cel{18}falsajo aspektar quale exaktajo, quale verajo. - DEF.}
\l{}
\l{}
\l{}
\l{\cel{44}------}
}
